\chapter{Formal Verification}\label{ch:formal}

Two TLA+ specifications were written for this project and are checked exhaustively by TLC on
every push. This chapter states what they establish, and — with equal emphasis — what they do
not, because the second specification is named \code{Isolation} and that name invites a stronger
reading than the model supports.

The chapter also records a defect in the first specification. For part of this project the
specification passed every check while being incapable of failing. That is reported here rather
than quietly repaired, because it is the most instructive methodological result in the work.

\section{Why formal methods were applied at all}

The firmware's safety claim is universal: \emph{no} sequence of adversarial messages causes an
unauthorised disarm. Testing cannot establish a universal claim over an unbounded input space;
it can only fail to refute it. The adversarial campaign of Chapter~\ref{ch:results} ran 9,600
trials and found no violation, which is evidence but not proof. Exhaustive model checking over
a bounded abstraction is the complementary instrument: it enumerates every reachable state of
the abstraction and therefore \emph{proves} the invariant within that abstraction's bounds.

The two instruments fail in different directions, which is why both are used. The model may
abstract away the very detail that breaks in reality; the experiment covers reality but only
the paths it happened to take.

\section{Specification 1: the alarm state machine}

\subsection{What is modelled}

\code{formal/AlarmFSM.tla} models the five-state arm/disarm machine and the authenticated
command path, over three variables:

\begin{itemize}[leftmargin=1.6em,itemsep=3pt]
  \item \code{state} --- one of \code{Disarmed}, \code{ExitDelay}, \code{Armed},
        \code{EntryDelay}, \code{Alarm}
  \item \code{ctr} --- the highest accepted authenticated command counter
  \item \code{badDisarm} --- set \code{TRUE} if an unauthenticated remote disarm ever takes
        effect
\end{itemize}

\textbf{Time is abstract.} A delay expires as a nondeterministic step rather than by counting
milliseconds. This is the correct level of abstraction for reasoning about \emph{which
transitions are reachable and who may cause them}, and it is the wrong level for reasoning about
\emph{how long} anything takes. No timing conclusion is drawn from this model anywhere in this
report.

\subsection{Properties checked}

\begin{longtable}{@{}p{0.27\textwidth}p{0.11\textwidth}p{0.50\textwidth}@{}}
\toprule
\textbf{Property} & \textbf{Kind} & \textbf{Statement} \\
\midrule
\endhead
\code{TypeOK}              & invariant & All three variables stay in their declared domains \\
\code{NoUnauthRemoteDisarm}& invariant & \code{badDisarm = FALSE}: the system never leaves an
armed state because of an unauthenticated remote command \\
\code{NoAlarmWhenDisarmed} & invariant & The siren is never on while disarmed \\
\code{CounterMonotonic}    & temporal  & $[][ctr' \geq ctr]$: the replay counter never
decreases \\
\code{EntryLeadsToResolve} & liveness  & An entry delay always ends in \code{Alarm} or
\code{Disarmed}; the system cannot get stuck mid-entry \\
\code{AlarmReArms}         & liveness  & The alarm always ends, in re-arm or disarm \\
\bottomrule
\end{longtable}

The liveness properties require fairness on the three timed transitions
(\code{ExitExpire}, \code{EntryExpire}, \code{AlarmExpire}), supplied by weak fairness in
\code{Spec}. No fairness is assumed for any adversary action.

\subsection{Result}

With \code{MaxCtr = 4} and \code{AllowUnauthDisarm = FALSE}: \textbf{166 states generated, 25
distinct, 0 left on queue, no error found.} The run is exhaustive over the bounded state space,
and both temporal branches were checked over the complete space.

\section{The vacuity defect, and the repair}\label{sec:formal-vacuity}

\subsection{What was wrong}

In the original specification the unauthenticated-command action was written as an
unconditional refusal:

\begin{lstlisting}[style=tla,caption={The defective action, as originally written (reconstructed
from the repair commit).},label={lst:vacuous}]
\* An UNauthenticated remote command is refused, exactly as the firmware does.
RemoteUnauthDisarm ==
  /\ state \in ArmedLike
  /\ UNCHANGED <<state, ctr, badDisarm>>
\end{lstlisting}

The comment asserts the conclusion. \code{badDisarm} was initialised \code{FALSE} and left
\code{UNCHANGED} by \emph{every} action in the specification, so no reachable behaviour could set
it. The invariant \code{NoUnauthRemoteDisarm} was therefore \textbf{unfalsifiable by
construction}: not proved, but incapable of failing.

TLC reported success. The success meant nothing.

\subsection{Why this mattered more than an ordinary bug}

The firmware contained a real defect of exactly the kind this invariant was written to catch: a
compile-time guard that admitted remote disarm without application-layer authentication. The
ablation experiment found it. The model checker, which existed specifically to find it, could
not have found it under any circumstances.

The specification had embedded the same unverified assumption that the firmware got wrong. A
passing formal check had been taken as evidence for a property that the check was structurally
unable to test. This is the project's central methodological finding, generalised in
Chapter~\ref{ch:conclusion}.

\subsection{The repair}

The build switch is now a model parameter rather than a hard-coded refusal:

\begin{lstlisting}[style=tla,caption={The repaired action: the configuration is a parameter, so
the violating behaviour exists and is reachable under one valuation.},label={lst:repaired}]
RemoteUnauthDisarm ==
  /\ state \in ArmedLike
  /\ IF AllowUnauthDisarm
       THEN /\ state' = "Disarmed"
            /\ badDisarm' = TRUE
            /\ UNCHANGED ctr
       ELSE UNCHANGED <<state, ctr, badDisarm>>
\end{lstlisting}

\subsection{The third configuration: a check that must fail}

The repair alone is not sufficient, because a future edit could reintroduce the vacuity and the
suite would again report success. A third configuration therefore exists whose \emph{required
outcome is a violation}.

\code{AlarmFSM\_unsafe.cfg} sets \code{AllowUnauthDisarm = TRUE}, modelling the
\code{ALLOW\_REMOTE\_DISARM} build with no application-layer authentication.
\code{formal/run\_tlc.sh} inverts the exit condition for this run and additionally requires the
specific string \code{Invariant NoUnauthRemoteDisarm is violated} in TLC's output:

\begin{lstlisting}[style=sh,caption={The inverted check. A \emph{passing} TLC run is treated as a
build failure.},label={lst:invertedcheck}]
if java -cp "$JAR" tlc2.TLC -config AlarmFSM_unsafe.cfg AlarmFSM.tla > /tmp/tlc_unsafe.out 2>&1
then
  echo "FAIL AlarmFSM_unsafe: TLC found no violation, so the invariant is vacuous again"
  fail=1
elif grep -q "Invariant NoUnauthRemoteDisarm is violated" /tmp/tlc_unsafe.out; then
  echo "PASS AlarmFSM_unsafe: TLC reports the expected NoUnauthRemoteDisarm violation"
\end{lstlisting}

Requiring the specific invariant name, rather than merely a non-zero exit, matters: a syntax
error, a missing constant or a parse failure also produce a non-zero exit, and would otherwise be
silently accepted as the expected violation.

TLC reports the expected violation on every run. The invariant is now known to be falsifiable,
which is what licenses reading its success in the safe configuration as evidence.

\section{Specification 2: the isolation model}\label{sec:formal-isolation}

\subsection{What is modelled}

\code{formal/Isolation.tla} models the two-task structure over seven variables:
\code{safetySteps}, \code{produced}, \code{queued}, \code{delivered}, \code{dropped},
\code{netState} and \code{netOwnsState}.

The network task is modelled \textbf{adversarially and with no fairness whatsoever}. It may move
between \code{idle}, \code{busy}, \code{stalled} and \code{down}; it may stall for an unbounded
time, as in a slow DNS, TCP or TLS call; it may fail to publish; and it may crash and restart.
Only the safety task is given weak fairness. The question the model answers is whether the
safety task still progresses when the network task is permitted to behave as badly as it likes
forever.

\subsection{Properties checked}

\begin{longtable}{@{}p{0.26\textwidth}p{0.11\textwidth}p{0.51\textwidth}@{}}
\toprule
\textbf{Property} & \textbf{Kind} & \textbf{Statement} \\
\midrule
\endhead
\code{TypeOK}               & invariant & All seven variables stay in their domains \\
\code{SafetyCanAlwaysStep}  & invariant &
$(\mathit{safetySteps} < \mathit{MaxEvents}) \Rightarrow \mathrm{ENABLED}\ \mathit{SafetyStep}$
--- in every reachable state the safety task is enabled \\
\code{NoNetworkWrite}       & invariant & \code{netOwnsState = FALSE}: the network task never
writes the alarm state \\
\code{QueueBounded}         & invariant & $\mathit{queued} \leq \mathit{QCap}$: the writer drops
instead of blocking \\
\code{NoSilentLoss}         & invariant &
$\mathit{produced} = \mathit{delivered} + \mathit{queued} + \mathit{dropped}$ \\
\code{SafetyProgress}       & liveness  &
$\Diamond(\mathit{safetySteps} = \mathit{MaxEvents})$: the safety task is never starved \\
\bottomrule
\end{longtable}

\code{SafetyCanAlwaysStep} is the architectural claim stated formally. Its force comes from what
the guard of \code{SafetyStep} mentions: \code{safetySteps} and \code{produced} only, and nothing
owned by the network side. Because the enabling condition is independent of \code{netState}, no
network behaviour can disable it. \code{NoSilentLoss} is the drop-oldest policy stated as an
accounting identity: an event is delivered, still queued, or counted as dropped, and never
disappears.

\subsection{Result}

With \code{QCap = 16} and \code{MaxEvents = 8}: \textbf{2,686 states generated, 660 distinct, 0
left on queue, no error found}, with the temporal property checked over the complete 660-state
space.

\subsection{What this model does and does not establish}

This distinction is the most important sentence in the chapter, and it is maintained verbatim
throughout this report and the manuscript.

\textbf{What it establishes.} \emph{Logical} isolation. The safety task's enabling condition is
independent of every network-side variable, so there exists no reachable state --- under
unbounded stalling, flooding, publication failure or crash --- in which the network side has
disabled the safety side. In the firmware this corresponds to the asymmetric mutex timeout: the
safety loop calls \code{xSemaphoreTake(statusMutex, 0)}, so its blocking term is identically
zero, by construction and not by measurement.

\textbf{What it does not establish.} \emph{Physical timing invariance.} The model has no notion
of microseconds, cache lines, bus arbitration, interrupt latency, radio duty cycle, or the
thermal and clock behaviour of a real die. Two cores on one package contend for shared silicon
in ways a seven-variable state machine cannot represent. Logical isolation does not imply that
the safety loop's measured wall-clock period is invariant to network load on real hardware, and
nothing in this project has measured whether it is.

The firmware is designed so that the \emph{blocking} component of that period is zero. Whether
the \emph{total} period is invariant is an empirical question about silicon, requiring hardware
conditions A--D. Those were not executed; Appendix~\ref{app:hardware} records why, and the
pre-registered analysis plan that was committed before any data existed.

\section{Trace conformance: connecting the model to the code}

A specification can be correct about a machine the firmware does not implement. Trace
conformance is the bridge.

\code{experiments/test\_conformance.py} parses the transition relation directly from
\code{AlarmFSM.tla} --- rather than re-encoding it in Python, where it could drift --- and checks
every recorded firmware state transition against it.

\textbf{19,529 transitions} across three ablation modes: 7,044, 5,282 and 7,206.
\textbf{Zero disallowed transitions.}

The figure is meaningful only because of the negative control. The same checker is given a
deliberately forbidden transition, \code{Disarmed} $\rightarrow$ \code{Alarm}, and is required to
reject it. It does. Without that control, a checker that accepted everything would also report
zero violations across 19,529 transitions, and the result would be indistinguishable from the
one reported here. This is the same reasoning as the inverted TLC configuration in
Section~\ref{sec:formal-vacuity}, applied to a different instrument.

\section{Summary}

\begin{longtable}{@{}p{0.23\textwidth}p{0.20\textwidth}p{0.45\textwidth}@{}}
\toprule
\textbf{Configuration} & \textbf{Result} & \textbf{Required outcome} \\
\midrule
\endhead
\code{AlarmFSM.cfg}        & 25 distinct states, no error  & pass \\
\code{Isolation.cfg}       & 660 distinct states, no error & pass \\
\code{AlarmFSM\_unsafe.cfg}& \code{NoUnauthRemoteDisarm} violated &
\textbf{violation required}; a pass is a build failure \\
\code{test\_conformance.py}& 19,529 transitions, 0 violations & pass, with the negative control
rejected \\
\bottomrule
\end{longtable}

All four are re-run by continuous integration on every push, so neither the vacuity defect nor a
silently accepting conformance checker can recur without the build failing.
